% ----------------------------------------------------------- %
% THIS IS WHERE YOU SHOULD ENTER ALL OF YOUR RELEVANT DETAILS %
% ----------------------------------------------------------- %

% Personal details

\renewcommand{\firstname}{Skye}
\renewcommand{\middlename}{Cassiopeia-Jacqueline}
\renewcommand{\lastname}{Owen-Lloyd-Walters}
\renewcommand{\postnomial}{MPhys MInstP FRAS}
\renewcommand{\pronoun}{She/They}
\renewcommand{\tagline}{Software \& Distributed Systems Engineer}
\renewcommand{\summary}{Software engineer specialising in distributed systems, platform engineering, and software supply-chain security. Sole owner of Canonical MAAS's CI platform and core maintainer of its Go Terraform provider; equally at home in Go services and Temporal orchestration or highly-available infrastructure-as-code. A First-Class MPhys astrophysicist who builds the systems other engineers rely on.}
\renewcommand{\profilephoto}{profile.jpg}
% contact
\renewcommand{\email}{skyecasolw@gmail.com}
\renewcommand{\address}{Havant, Hampshire, United Kingdom}

% social details

\renewcommand{\website}{sk1y101.github.io}
\renewcommand{\githubusername}{sk1y101}
\renewcommand{\linkedinusername}{skye-owen-lloyd-walters}
\renewcommand{\orcidid}{0000-0003-2186-1582}

% Colours and settings
% red, orange, green, aquamarine, aqua, blue, lilac, purple, dark grey, grey, light grey, lightergrey

\colorlet{titlecolour}{lightgrey}
\colorlet{backgroundcolour}{white}
\colorlet{sidebarcolour}{lightergrey}
\colorlet{highlightcolour}{lilac}
\colorlet{skillcolour}{lilac}
\renewcommand{\slimdates}{true}

% Languages
% \languageskill{language code}{proficiency}
% Proficiency is rated: 1-5, where 1 is beginner, 5 is fluency. -1 is additionally used for a native language

\languageskill{English}{-1}
\languageskill{Japanese}{1}
% \languageskill{German}{1}

% Skills; unified command {category}{name}{level}
% Category used for grouping in resume.json; all skills shown together in sidebar
% Mastery: 5=Expert/Senior, 4=Advanced, 3=Intermediate, 2=Mediocre, 1=Beginner

\iflean\else
\skill{Physics}{Computational Modelling}{4}
\skill{Physics}{Signal Processing}{4}
\skill{Physics}{Statistical Analysis}{4}
\fi

\skill{Languages}{Bash}{3}
\skill{Languages}{Go}{4}
\skill{Languages}{Groovy}{4}
\skill{Languages}{Javascript}{3}
\skill{Languages}{Python}{5}
\skill{Languages}{SQL}{3}
\iflean\else
\skill{Languages}{HTML/CSS}{3}
\skill{Languages}{\LaTeX}{4}
\skill{Languages}{Markdown}{4}
\fi

\skill{DevOps \& Cloud}{Ansible}{4}
\skill{DevOps \& Cloud}{Docker}{3}
\skill{DevOps \& Cloud}{Git}{4}
\skill{DevOps \& Cloud}{Jenkins / CI}{5}
\skill{DevOps \& Cloud}{Juju / Charmed Operators}{3}
\skill{DevOps \& Cloud}{LXD}{4}
\skill{DevOps \& Cloud}{Linux}{4}
\skill{DevOps \& Cloud}{Networking}{4}
\skill{DevOps \& Cloud}{PostgreSQL}{3}
\skill{DevOps \& Cloud}{REST APIs}{4}
\skill{DevOps \& Cloud}{Terraform HCL}{4}
\skill{DevOps \& Cloud}{Test Automation}{4}
\skill{DevOps \& Cloud}{YAML}{3}

\skill{Backend \& Distributed Systems}{Distributed Systems}{4}
\skill{Backend \& Distributed Systems}{Temporal}{4}
\skill{Backend \& Distributed Systems}{Event-Driven Architecture}{3}
\skill{Backend \& Distributed Systems}{High Availability}{4}
\skill{Backend \& Distributed Systems}{Supply-Chain Security (SBOM/CVE)}{4}

\iflean\else
\skill{Web \& Data}{Jupyter}{3}
\skill{Web \& Data}{NumPy / SciPy}{5}
\fi

% Jobs
% \addcompany{Company name}{logo path}
% \addposition{Company name}{Role}{team}{Start}{end}{details}

\addcompany{Canonical}{canonical_logo.jpeg}
\addposition{Canonical}
{Software Engineer}{Canonical MAAS (Metal-As-A-Service)}
{2023-05}{}
{%
\begin{itemize}[leftmargin=*, labelsep=-0.1em, noitemsep]
    \item \textbf{Sole owner and architect} of the MAAS team's entire Jenkins CI platform; \textbf{130+ pipeline jobs} generated from \textbf{39 reusable templates} across \textbf{42 product repositories}, defined as code in $\sim$17k lines of Python, Groovy, and YAML. Designed the templated job system, automated bug backporting from merge proposals (lander/reviewer polling every two minutes), drove the Launchpad-to-GitHub migration, and introduced testflinger-driven bare-metal hardware testing.
    \item Architected a \textbf{Temporal-based distributed orchestration system} that builds, tests, and promotes MAAS OS images at scale, engineering production-grade resilience (activity heartbeats, execution timeouts, orphaned-workflow termination, and worker drain/health monitoring) and a custom workflow scheduler to replace unreliable SDK scheduling.
    \item Built MAAS's \textbf{software supply-chain security pipeline} from scratch: a 1,360-line CVE/SBOM scanner covering snap and deb packages across \textbf{5 MAAS release lines}, integrating Trivy, BlackDuck, Canonical's SBOM service, and the Ubuntu CVE tracker, with \textbf{dynamic SBOM generation} via resumable chunked uploads.
    \item \textbf{Core maintainer} of the \href{https://github.com/canonical/terraform-provider-maas}{MAAS Terraform Provider} ($\sim$20k LOC Go), \textbf{authoring} its storage subsystem (software RAID, LVM volume groups and logical volumes) and boot-image streaming resources end-to-end, with full CRUD lifecycle, schema validation, and state migration.
    \item Designed \textbf{highly-available} Charmed MAAS deployments in \href{https://github.com/canonical/maas-terraform-modules}{Terraform/Terragrunt} (three-node region, database, and proxy topologies with HAProxy + Keepalived virtual-IP failover) and implemented cluster-coordinated, S3-backed backup and restore in the \href{https://github.com/canonical/maas-charms}{MAAS region charm}.
    \item Contributed idiomatic Go to the \href{https://github.com/canonical/gomaasclient}{gomaasclient} API client; interface-segregated endpoint design and defensive deserialization (multi-format timestamp parsing, numeric-type coercion) against a live REST API.
\end{itemize}
}

\addposition{Canonical}
{Associate Software Engineer}{Canonical MAAS (Metal-As-A-Service)}
{2022-06}{2023-05}
{%
\begin{itemize}[leftmargin=*, labelsep=-0.1em, noitemsep]
    \item Selected for the MAAS team through Canonical's competitive \textbf{graduate fast-track} as a Python backend developer.
    \item \textbf{Authored} the \href{https://github.com/canonical/ansible-collection}{MAAS Ansible Collection}, published to \textbf{Ansible Galaxy}, giving operators declarative automation of MAAS environments, alongside the underlying \href{https://github.com/maas/MAAS-ansible-playbook}{playbooks} for bootstrapping and configuring deployments.
    \item \textbf{Authored} the \href{https://maas.io/docs/api}{MAAS OpenAPI schema}, documenting the entire MAAS REST API for language-agnostic client generation; the foundation later consumed by the Go provider and client I went on to maintain.
    \item \textbf{Initiated} the \href{https://github.com/canonical/terraform-provider-maas}{MAAS Terraform Provider} with OpenAPI-driven resource management; work I would later lead as core maintainer.
    \item Built \href{https://github.com/canonical/packer-maas}{Packer templates} producing custom deployable OS images for MAAS environments.
\end{itemize}
}

% Education
% \addeducation{School name}{logo path}
% \addcourse{School name}{Course name}{final grade}{Start}{end}{details}

\addeducation{University of Portsmouth}{UOP_logo.png}
\addcourse{University of Portsmouth}
{MPhys Physics, Astronomy, and Cosmology}{First Class Honours}
{2018-09}{2022-05}
{%
\begin{itemize}[leftmargin=*, labelsep=-0.1em, noitemsep]
    \item First year mean grade of 76.7\%, Final year mean grade of 79.2\%
    \item \href{https://sk1y101.github.io/assets/pdf/TransitDissertation.pdf}{Masters thesis} combining computational modelling, first hand observation, and supplemental data to measure exoplanet transits. Developed analytical models of transit timing variation to investigate the orbital properties of non-transiting planets. Results were summarised with a 6000 word dissertation, 10 minute presentation/discussion, and scientific poster.
    \item \href{https://sk1y101.github.io/assets/pdf/GWProjectFinalReport.pdf}{Bachelors thesis} utilising computational modelling and signal processing to identify and model glitch events within the LIGO data set. Results were summarised with a 5000 word dissertation, 10 minute presentation/discussion, and LaTeX Compatible result list.
\end{itemize}
}

\addeducation{AiCore}{aicore_logo.png}
\addcourse{AiCore}
{Ai and Data Engineering}{Certified in the practical application of AI and Data Engineering}
{2022-01}{2022-06}
{%
\begin{itemize}[leftmargin=*, labelsep=-0.1em, noitemsep]
    \item Completed the course to the highest standard awarded in less than half the expected time.
\iflean
    \item Software, data, and cloud engineering: advanced Python, SQL, Docker, and the AWS serverless stack.
\else
    \item Software engineering: Git \& GitHub, advanced Python, algorithms \& data structures.
    \item Data engineering: SQL, data lakes, data warehousing, web scraping with Selenium.
    \item Data science: Data cleaning, preprocessing \& visualisation, A/B testing, feature engineering, statistical modelling, model selection and implementation using sklearn.
    \item Cloud engineering: cloud computing, designing and building APIs, Docker, Apache Airflow, AWS Serverless Stack (EC2, S3, RDS).
\fi
\end{itemize}
}

\iflean\else
\addeducation{Peter Symonds College}{petersymonds_logo.jpeg}
\addcourse{Peter Symonds College}
{A-Level}{Physics (A), EPQ (A), Maths (B), Further Maths (D), AS-Chemistry (A)}
{2016-09}{2018-06}
{%
\begin{itemize}[leftmargin=*, labelsep=-0.1em, noitemsep]
    \item Produced a Level 3 Extended Project Qualification to a high standard titled \href{https://sk1y101.github.io/assets/pdf/EPQ.pdf}{`Is There A Possibility of Extrasolar Habitation?'}, assessing the habitability of the TRAPPIST-1 system and its seven Earth-sized exoplanets orbiting an M8V ultracool dwarf star.
\end{itemize}
}

\addeducation{Trafalgar School}{Trafalgar_school_logo.png}
\addcourse{Trafalgar School}
{GCSE}{9 GCSE's A* to B Including Science, Computer Science, Maths, and English}
{2011-09}{2016-05}
{}
\fi

% Professional Memberships
% \addmembership{Location}{logo path}
% \addmembershiptitle{Location}{Title}{onelinedetauks}{Start}{end}{details}

\iflean\else
\addmembership{British Astronomical Association}{baa_logo.jpeg}
\addmembershiptitle{British Astronomical Association}
{Member}{Exoplanet transit timing and observation}
{2022}{}
{}
{https://britastro.org/}
\fi

\addmembership{Institute of Physics}{IOP_logo.png}
\addmembershiptitle{Institute of Physics}
{Member}{Granted the use of the post-nomen `MInstP'}
{2022}{}
{}
{https://www.iop.org/}
\iflean\else
\addmembershiptitle{Institute of Physics}
{Associate Member}{Affiliate membership during MPhys studies}
{2018}{2022}
{}
{https://www.iop.org/}

\addmembership{European Astronomical Society}{eas_logo.jpeg}
\addmembershiptitle{European Astronomical Society}
{Member}{European astronomy and astrophysics research}
{2020}{}
{}
{https://eas.unige.ch/}
\fi

\addmembership{Royal Astronomical Society}{ras_logo.jpeg}
\addmembershiptitle{Royal Astronomical Society}
{Fellow}{Granted the use of post-nomen `FRAS'}
{2020}{}
{}
{https://ras.ac.uk/}

% Awards
% \addaward{award name}{affiliated location}{date received}{details}

\addaward{Certified in the practical application of AI and Data Engineering}
{AiCore}{2022-06}
{Awarded for the successful completion of the course. Completed to the highest standard awarded in less than half the expected time.}

\addaward{Graham Bryant Prize for Best Observatory Project}
{University of Portsmouth}{2022-05}
{First recipient of the award to honour the late Graham Bryant for my Masters thesis}

% Projects and Publications

% \addproject{name}{icon}{affiliation}{date}{details}{link}

\addproject{Curriculum Vitae; LaTeX}{code}
{Personal}{2026-05-18--Present}
{\iflean An AI-native CV build system: a custom LaTeX document class that renders both this lean CV and an exhaustive website version from one source, with a Python JSON-Resume exporter and LLM-generated project summaries (model fallback, rate-limiting, health tracking), released via CI on every push.\else A fully custom LaTeX document class and build system that renders two CVs (this lean job version and an exhaustive website version) from a single source via a lean/full switch. A Python toolchain exports the JSON Resume schema and synchronises GitHub repositories into project entries, generating their descriptions through an LLM with model fallback, request rate-limiting, and health tracking. GitHub Actions builds and releases both PDFs and the JSON on every push to main.\fi}
{https://github.com/SK1Y101/cv}

\addproject{MAAS Internal CI}{code}
{Canonical}{2023-03--Present}
{\iflean The MAAS team's Jenkins CI platform (Python, Groovy, YAML); CI/CD pipelines, supply-chain security (CVE and SBOM) scanning, and release automation across 42 product repositories.\else Sole architect and owner of the MAAS team's entire Jenkins CI platform; 130+ pipeline jobs generated from 39 reusable templates across 42 product repositories, defined as code in roughly 17,000 lines of Python, Groovy, and YAML. Built the team's software supply-chain security pipeline from scratch: a 1,360-line CVE and SBOM scanner covering snap and deb packages across five MAAS release lines, integrating Trivy, BlackDuck, Canonical's SBOM service, and the Ubuntu CVE tracker. Designed the reusable templated job system (a single config file provisions coordinated lander, reviewer, and tester jobs that poll every two minutes), automated bug backporting from merge proposals, drove the Launchpad-to-GitHub migration, introduced testflinger-driven bare-metal hardware testing, and built a live CI observability dashboard with Temporal-backed failure tracking.\fi}
{https://maas.io}

\addproject{MAAS Terraform Provider}{code}
{Canonical}{2023-08--Present}
{\iflean Canonical's official Go Terraform provider for MAAS (~20k LOC, HashiCorp Plugin SDK v2); core maintainer, authoring the storage and boot-image streaming resources.\else Core maintainer of Canonical's official Terraform provider for MAAS ($\sim$20,000 lines of Go on the HashiCorp Plugin SDK v2). Authored the storage subsystem (software RAID, LVM volume groups and logical volumes) and the boot-image streaming resources end-to-end, with full CRUD lifecycle, schema validation, and state migration; and owns the provider's Go CI and release pipeline.\fi}
{https://github.com/canonical/terraform-provider-maas}

\addproject{MAAS System Tests}{code}
{Canonical}{2022-09--Present}
{\iflean MAAS's end-to-end test platform; a Temporal-based distributed pipeline that builds, tests, and promotes OS images at scale across VMs and real bare-metal hardware.\else Principal engineer behind MAAS's end-to-end test platform. Architected a Temporal-based distributed orchestration system that builds, tests, and promotes OS images at scale across LXD, QEMU, and real bare-metal hardware, engineering production-grade resilience: a custom workflow scheduler, activity heartbeats, execution timeouts, orphaned-workflow termination, and worker drain and health monitoring.\fi}
{https://maas.io}

\addproject{MAAS Terraform Modules}{code}
{Canonical}{2025-09--Present}
{\iflean Terraform/Terragrunt infrastructure-as-code deploying highly-available Charmed MAAS with HAProxy and Keepalived virtual-IP failover.\else Lead author of the infrastructure-as-code that deploys highly-available Charmed MAAS. Designed three-node region, database, and proxy topologies with HAProxy and Keepalived virtual-IP failover, S3 backup integration, and dependency-ordered Terragrunt stacks; reducing a full HA cluster deployment to a few commands.\fi}
{https://github.com/canonical/maas-terraform-modules}

\addproject{MAAS Charms}{code}
{Canonical}{2025-07--Present}
{\iflean Juju charmed operators that run MAAS; built the cluster-coordinated, S3-backed backup and restore subsystem and MAAS HAProxy TCP routing.\else Significant contributor to the Juju charmed operators that run MAAS. Built the cluster-coordinated, S3-backed backup and restore subsystem (streamed tar archiving, version verification, cross-charm coordination) and implemented MAAS HAProxy TCP routing alongside the region-and-rack charm consolidation.\fi}
{https://github.com/canonical/maas-charms}

\addproject{gomaasclient}{code}
{Canonical}{2024-06--2025-05}
{\iflean Go client library for the MAAS REST API, with an interface-segregated api/entity/client architecture; powers the Terraform provider.\else Contributor to the Go client library underpinning MAAS tooling and the Terraform provider. Added typed CRUD resources across an interface-segregated api/entity/client architecture and hardened deserialization against a live REST API; multi-format timestamp parsing, numeric-type coercion, and snake\_case operation naming.\fi}
{https://github.com/canonical/gomaasclient}

\addproject{Skiylia}{code}
{Personal}{2022-12-11--Present}
{\iflean A dynamically-typed, object-oriented programming language of my own design, compiling to bytecode for a custom stack-based virtual machine. Built in Python end-to-end; lexer, parser, bytecode compiler, and runtime; as a deep dive into language and compiler construction, documented on Read the Docs.\else A dynamically-typed, object-oriented programming language of my own design that compiles to bytecode for a custom virtual machine. Implemented in Python as a deep dive into compiler construction: lexing, parsing, bytecode compilation, and runtime execution. Documented on Read the Docs with Snyk vulnerability scanning and automated releases.\fi}
{https://github.com/SK1Y101/Skiylia}

\addproject{YouTube Automation Suite}{code}
{SkiylianSoftware}{2025-02-02--Present}
{\iflean A Python automation suite running my YouTube channel: OAuth integration with the YouTube and Google Calendar APIs to auto-sort videos into playlists, schedule released and upcoming videos across calendars, and assemble Shotcut video projects with background music; with a re-authentication flow and persisted credentials.\else A Python automation suite that runs my YouTube channel end-to-end: OAuth integration with the YouTube and Google Calendar APIs to auto-sort videos into playlists, schedule released and upcoming videos across calendars, and auto-populate Shotcut project files with background music; event-driven and credential-persisting.\fi}
{https://github.com/SkiylianSoftware/YT-automation}

\addproject{Exoclock Data Release IV}{pub}
{ExoClock}{2025-11}
{Peer-reviewed paper in the ExoClock collaboration's Data Release IV, presenting a homogeneous catalogue of 620 updated exoplanet ephemerides. Contributed observational data and transit timing analysis from the TransitProject pipeline. Published on arXiv.}
{https://arxiv.org/abs/2511.14407}

\iflean\else
\addproject{`Variations on an Exoplanet Theme' Webinar}{talk}
{British Astronomical Association}{2023-09}
{Computational Modelling of Transit Timing Variations}
{https://britastro.org/event/exoplanet-2023sep}

\addproject{Exoplanet TTV Masters Thesis}{pub}
{University of Portsmouth}{2022-05-25}
{Master's thesis (`Determining The Parameters of Exoplanetary Candidates From Transit Timing Variations') developing analytical Transit Timing Variation models to probe the dynamical configuration of exoplanetary systems and recover the parameters of non-transiting companion planets. Collected new transit light curves with the 24-inch Ritchey-Chr\'etien telescope at Clanfield Observatory and combined them with the ExoClock and Exoplanet Transit databases, then built a computational pipeline that fits competing TTV models and selects between them via the Akaike Information Criterion. Observations were contributed to the ExoClock collaboration's updated ephemerides catalogue.}
{https://sk1y101.github.io/assets/pdf/TransitDissertation.pdf}

\addproject{IMBH Glitch Detection Bachelors Thesis}{pub}
{University of Portsmouth}{2021-05}
{Bachelor's thesis (`Distinguishing Intermediate Mass Black Hole Mergers From Short Duration Glitches') building a computational model of the transient glitches that contaminate LIGO data at a rate of order ten per hour and mimic intermediate-mass black hole mergers. Validated the model against real LIGO strain data using matched filtering and cross-checked it against events flagged by the Omicron scan, demonstrating strong agreement despite the differing detection methods. Showed the glitch model does not suppress genuine signals, with the GW190521 merger templates still responding more strongly than the glitch template.}
{https://sk1y101.github.io/assets/pdf/GWProjectFinalReport.pdf}

\addproject{Extrasolar Habitation EPQ}{pub}
{Peter Symonds College}{2018-06}
{Level 3 Extended Project Qualification (`Is There A Possibility of Extrasolar Habitation?', graded A) evaluating the habitability of the TRAPPIST-1 system and its seven Earth-sized planets orbiting an M8V ultracool dwarf star. Weighed each world against habitability benchmarks spanning the circumstellar habitable zone, effective and atmospheric temperature, and long-term orbital stability, and used transit photometry to place upper bounds on the size of any potential exomoons. Concluded that the system most likely underwent planetary migration to reach its present configuration.}
{https://sk1y101.github.io/assets/pdf/EPQ.pdf}
\fi

% hobbies and interests
\addhobby{Gaming}{\iflean Factorio, Stellaris, Transport Fever, and Kerbal Space Program; there's something satisfying about optimising a sprawling factory or planning a multi-planetary slingshot.\else I Love Factorio, Stellaris, Transport Fever, modded Minecraft, and Kerbal Space Program; the latter three feature heavily in my YouTube content creation. There's something satisfying about optimising a factory or planning and executing a multi-planetary slingshot.\fi}
\addhobby{Programming}{\iflean I build Fitbit watch faces, design my own programming languages, write flight computers for Kerbal Space Program, and tinker with esoteric hardware like balanced-ternary logic.\else I build watch faces for my FitBit Versa, design programming languages, write flight computers for Kerbal Space Program, and experiment with esoteric hardware designs like balanced ternary logic. The thing I'm most proud of is the shotcut automation pipeline that backs my YouTube workflow.\fi}
\addhobby{Amateur Astronomy}{\iflean Exoplanet transit timing and observation through the ExoClock collaboration; catching a tiny dip in a star's brightness from my back garden never gets old.\else Exoplanet transit timing and observation through the ExoClock collaboration. There's something magical about catching a tiny dip in a star's brightness from my back garden, knowing I'm watching another world pass by.\fi}
\addhobby{Writing}{\iflean Documentation, the occasional blog post, and hobbyist short fiction; I find that explaining something clearly is the best way to understand it myself.\else Be it Documentation or the occasional blog post. I find that explaining something clearly is the best way to make sure I understand it myself. But it's hobbyist writing too, I enjoy partaking in short novella and building out the world that lives within my mind.\fi}
\iflean\else
\addhobby{Space \& Spaceflight}{Natural or artificial, if it's found outside of Earth's atmosphere, I probably like it. Exoplanets and deep space probes are my favourite. I keep an eye on every mission that leaves Earth.}
\addhobby{Reading}{I love a good book, but really wish I took the time to read more. My current favourites are Project Hail Mary and The Hydrogen Sonata; I have a soft spot for first-contact and post-scarcity sci-fi that explores how we might grow as a species.}
\addhobby{Language Learning}{I'm learning Japanese. Currently at a beginner level and making slow but steady progress; it's a beautifully efficient language.}
\addhobby{Tea}{Earl Grey is my favourite; no sugar, no milk. I do love a nice loose leaf brew from Whittards though, especially something fruity, and always stewed for just the perfect amount of time.}
\addhobby{Board Games}{My current top three are Hex Effects, Pandemic, and Mage Noir. I'm not great at strategy or deck-building, but I love playing them nonetheless. If it's a Co-op boardgame, I'll always have a blast.}
\fi
